% =============================================================================
% Results sec 5.5 -- Learning-rate asymmetry, FedNova regime-dependence
% Compressed version. Target ~560 words.
% =============================================================================

\section{Learning-rate asymmetry reveals FedNova regime-dependence}
\label{sec:lr-asymmetry}

On L4, Client 0 (the dominant client, $\approx 90\,\%$ of training samples)
keeps $\textrm{lr} = 0.01$ throughout while Client 1 (rare classes) takes
one of six learning-rate values, defining six asymmetry ratios from 1:1 to
50:1. Each cell is run across FedAvg, FedProx, and FedNova ($n=3$ matched
seeds per cell, $54$ runs total). Test macro-F1 across the envelope splits
the three optimisers cleanly:

\begin{center}
\begin{tabular}{c|cc|ccc|cc}
\hline
Ratio & Client 0 lr & Client 1 lr & FedAvg & FedProx & FedNova & FN$-$FA & FN$-$FP \\
\hline
 1:1  & $0.01$ & $0.01$    & $0.457$ & $0.491$ & $\mathbf{0.509}$ & $+0.052$ & $+0.018$ \\
 2:1  & $0.01$ & $0.005$   & $0.425$ & $0.462$ & $\mathbf{0.526}$ & $+0.101$ & $+0.064$ \\
 5:1  & $0.01$ & $0.002$   & $0.375$ & $0.392$ & $\mathbf{0.504}$ & $+0.129$ & $+0.112$ \\
10:1  & $0.01$ & $0.001$   & $0.369$ & $0.369$ & $\mathbf{0.501}$ & $+0.132$ & $+0.132$ \\
20:1  & $0.01$ & $0.0005$  & $0.367$ & $0.372$ & $\mathbf{0.486}$ & $+0.119$ & $+0.114$ \\
50:1  & $0.01$ & $0.0002$  & $0.367$ & $0.367$ & $\mathbf{0.498}$ & $\mathbf{+0.131}$ & $\mathbf{+0.131}$ \\
\hline
\end{tabular}
\end{center}

FedAvg and FedProx drop $\approx 0.10$ macro-F1 between 1:1 and 5:1 and then
plateau at the rare-class floor near $0.367$; FedNova stays in the
$0.486$--$0.526$ band across the entire envelope ($|\Delta$ vs its own 1:1
baseline$|$ $\leq 0.024$). At 50:1 the FedNova--FedAvg gap is $+0.131$ macro-F1
--- the largest single-axis system-asymmetry gap in this thesis and an order
of magnitude larger than the engineered statistical-heterogeneity
$\Delta = +0.007$ in \S\,\ref{sec:statistical-heterogeneity}. Rare-class F1
(mean across dermato, melanoma, vascular) collapses for FedAvg and FedProx
(FA $0.016$, FP $0.010$ at 10:1; $\approx 0$ by 50:1) while FedNova holds
$0.33$--$0.41$ throughout (FN $0.366$ at 10:1, $0.353$ at 50:1).

A mechanism-consistent reading: scaling Client 1's $\textrm{lr}$ by $1/5$
shrinks its per-round update magnitude correspondingly, and sample-weighted
aggregation does not normalise this away. FedNova's $\tau$-normalisation
divides each client's update by its effective work, which would restore
approximate contribution parity in this regime. The empirical update-norm
trace is consistent with this reading: at the 5:1 cell, Client 1's mean
update-norm ratio (5:1 $\div$ 1:1) is $0.47$ for FedAvg, $0.50$ for FedProx,
and $0.91$ for FedNova. This is mechanism-consistent evidence, not a
mechanism proof.

The same FedNova that preserves macro-F1 under deterministic LR asymmetry
fails under random-$\tau$ straggler heterogeneity. On
\texttt{system\_het\_random\_fednova} (L4, per-round random straggler
schedule, $\gamma$-inexact, $n=10$), mean macro-F1 is $0.237 \pm 0.094$, far
below FedAvg and FedProx on the matching \texttt{system\_het\_random} cells
(FA $0.493$, FP $0.510$ from \S\,\ref{sec:li-straggler}). Three of ten seeds
(s31337, s271828, s161803) collapse to macro-F1 $= 0.114$ --- the majority-
class-only baseline (predicting mel-nevi on every test image) at best-validation
selection; by the final validation round all ten runs fall below $0.20$ macro-F1,
nine reaching the strict mel-nevi-only pattern (final \emph{test} metrics were not
saved for this sweep, so this is a validation-trajectory observation). This regime
lies outside Wang 2020's stable-$\tau$ assumption: $\tau$-normalisation
amplifies stragglers' noisy partial updates rather than correcting for them.
FedNova's robustness is regime-dependent \emph{in both directions} ---
strongly robust under persistent LR asymmetry, brittle under random-$\tau$
instability.

\begin{center}
\begin{tabular}{lll}
\hline
FedNova regime              & Result                                   & Interpretation \\
\hline
Persistent LR asymmetry     & macro-F1 holds $0.49$--$0.53$ across 1:1--50:1 & Empirically robust; normalised aggregation, not direct LR correction \\
Random-$\tau$ stragglers    & macro-F1 = $0.237 \pm 0.094$; 3/10 collapse at best-val, 9/10 by final round & $\tau$-normalisation amplifies noisy partial updates (outside Wang 2020) \\
\hline
\end{tabular}
\end{center}

The escalation IID $\to$ symmetric L4 $\to$ Li 2020 contrast $\to$
perfect-storm $\mu = 1.0$ $\to$ perfect-storm $\mu = 0.01$ has within-pair
$\Delta$ rising from $-0.007$ and $+0.005$ (inside seed noise) to $+0.115$,
$+0.278$, $+0.404$ (large, mechanism-isolable, regime-driven). The
optimiser ranking is a function of which regime cell the protocol occupies,
not a property of the optimiser in the abstract.

The weak protocols fail in a clinically specific way: they lose rare-class
sensitivity, detailed next.


% =============================================================================
% Figure 5.5 placeholder
% =============================================================================

\begin{figure}[h]
\centering
\includegraphics[width=\textwidth]{F_asymmetric_lr_L4.pdf}
\caption{Learning-rate asymmetry envelope at L4, $n=3$ per cell. Two panels:
test macro-F1 and rare-class F1 (mean over dermato, melanoma, vascular) vs
the Client 1 LR ratio on $\log_{10}$ axis from 1:1 to 50:1. FedAvg and
FedProx plateau at the rare-class floor; FedNova stays flat across the
envelope.}
\label{fig:lr-asymmetry-envelope}
\end{figure}



